\section{External validation and operational study protocol}\label{sec:validation}
\subsection{Minimum admissible export}
An independent evaluation needs more than alert labels. A SOC partner must supply a versioned source-observation identifier, trustworthy upstream-root and derivation ancestry, evidence-arrival time, event time, causal process/authentication/network fields, organization and queue membership, and human-adjudicated alert-to-incident assignments. Priority judgments must specify the candidate universe, ordering direction, completeness, query cutoff and how ties or unjudged incidents are handled. These are preconditions for reproducing the frozen decision at the time it would have been made. The available GUIDE fields do not satisfy them; imputing distinct roots from shared host/account values or substituting noised creation timestamps for availability would test a different algorithm.

Before model execution, a data steward should approve field purpose and visibility, resolve or explicitly exclude ambiguous alert membership under a frozen policy, and partition organizations or time windows without post-outcome leakage. The first qualifying cohort should be used to verify schema and calibrate operational thresholds only; a disjoint development cohort should select candidates; an untouched test cohort should be sealed before any label access. Every arm must receive identical eligible alerts, candidate pairs, visible evidence and cutoff times. Group caps and abstention behavior must be fixed in advance. Hashes of code, configuration, inputs, predictions and decision traces provide a checkable chain from request to score.

\subsection{Comparison and analysis plan}
Prespecify a rule-based clusterer, a lexical retrieval control, a single-call LLM, a repeated single-LLM control with equal total provider budget, C+P without V, full C/P/V, and the deterministic verifier. The measured synthetic holdout comparison in this paper is \emph{only} the last method against its frozen matched baseline; it cannot fill this future comparison table. For a fair LLM experiment, report model ID/version, prompts, decoding, actual input/output tokens, retry counts, latency, cost, invalid JSON, citation validity, review coverage, and how a timed-out request contributes to denominators. The same candidate universe must be offered to every method; otherwise candidate recall changes are confounded with agent reasoning.

Primary endpoints should be incident Recall@20 under the predeclared 80\% completeness/purity rule and all-group false grouped pairs. Secondary endpoints include P@10, nDCG@20, recoverability, candidate coverage, pair precision/recall with explicit eligible-pair universe, reviewer citation agreement, review burden and latency quantiles. Pair-level TN requires an exhaustive specified pair universe; a positive-only expert-priority list cannot yield classification accuracy. Report per-queue paired differences and confidence intervals using queue or organization as the independent unit, and state multiplicity control for multiple comparisons. A reviewer study should measure time to first correct triage, workload and harmful false merges under human supervision. No automated containment or incident closure follows from the offline evidence.

\subsection{Decision gates and reporting}
External promotion should require a predeclared non-inferiority bound for false merges and useful retrieval, acceptable review/abstention rates, reliable point-in-time operation, and a trace review by analysts. A zero-error synthetic observation does not set those bounds. Document common-mode sensor failures, correlated roots, spoofed fields, clock drift, late arrivals, organization leakage and large incidents as stress strata. If a qualification or authorization gate fails, report ``not evaluated'' and retain the previous frozen method; do not silently relax the verifier to create a score. The present real-data audit is such a failure of test readiness, and production/shadow deployment is therefore not an outcome of this study.
